\begin{proof}[Proof of \cref{obj:thm:gap-window}]
1.  \Cref{obj:ass:two-block-homophily} gives \(m\ge2\), \(b<a\), and \(b>0\). Hence
\[
q=2(m-1)>0,\qquad \sqrt q>0,\qquad a+b>0,\qquad a-b>0.
\]
\cref{obj:ass:low-scale-two-block} gives
\[
2m-a-3b>0.
\]
These inequalities are the standing sign facts used below for division and multiplication by positive quantities.  % lean: gap_window

2.  Assume now that \cref{obj:ass:odd-community-size} holds, and fix
\[
0\le \kappa<\kappa_{\mathrm{gap}}(m,a,b)=\frac{(2m-a-3b)(a-b)\sqrt q}{a+b}.
\]
Multiplying this strict inequality by the positive factor \((a+b)/((a-b)\sqrt q)\),
\[
\frac{\kappa(a+b)}{(a-b)\sqrt q}<2m-a-3b.
\]
Next, from \(1+\dfrac{2b}{a-b}=\dfrac{a+b}{a-b}\) we get the identity
\[
R_x^+-R_x^-
=(a+b)\Bigl(2m-a-b-\frac{\kappa}{\sqrt q}\Bigr)-2b(a+b)\Bigl(1+\frac{\kappa}{(a-b)\sqrt q}\Bigr)
=(a+b)\left((2m-a-3b)-\frac{\kappa(a+b)}{(a-b)\sqrt q}\right),
\]
whose right-hand side is positive by the previous display and \(a+b>0\).  Thus \(R_x^-<R_x^+=r_{\mathrm{gap}}^+\), and since
\[
r_{\mathrm{gap}}^-=\frac{R_x^-+R_x^+}{2}
\]
is the midpoint of \(R_x^-\) and \(R_x^+\),
\[
R_x^-<r_{\mathrm{gap}}^-<r_{\mathrm{gap}}^+.
\]
Also \(R_x^-\ge0\), because \(2b(a+b)>0\) and \(\kappa/((a-b)\sqrt q)\ge0\); and since \(1-\kappa/(a-b)\le1\le1+\kappa/((a-b)\sqrt q)\) and \(2b(a+b)>0\),
\[
2b(a+b)\left(1-\frac{\kappa}{a-b}\right)\le R_x^-,
\]
so, the maximum of two quantities each at most \(R_x^-\) being at most \(R_x^-\),
\[
r_{\mathrm{cut,gap}}(m,a,b,\kappa)
=\max\left\{0,\;2b(a+b)\left(1-\frac{\kappa}{a-b}\right)\right\}
\le R_x^-.
\]
Combining these inequalities yields
\[
r_{\mathrm{cut,gap}}(m,a,b,\kappa)
<
r_{\mathrm{gap}}^-(m,a,b,\kappa)
<
r_{\mathrm{gap}}^+(m,a,b,\kappa).
\]  % lean: gap_window

3.  Let \(r_{\mathrm{gap}}^-<r<r_{\mathrm{gap}}^+\). From the preceding step,
\[
R_x^-<r.
\]
After division by \(2b(a+b)>0\), this says
\[
1+\frac{\kappa}{(a-b)\sqrt q}<\frac{r}{2b(a+b)}.
\]
Multiplying by \(a-b>0\) gives
\[
\frac{\kappa}{\sqrt q}
<
(a-b)\left(\frac{r}{2b(a+b)}-1\right).
\]
Using the identity
\[
2b+\frac r{2b}-\left(a+b+\frac r{a+b}\right)
=
(a-b)\left(\frac{r}{2b(a+b)}-1\right),
\]
and \(c_x/q=a+b+r/(a+b)\), this proves
\[
c_y>\frac{c_x}{q}+\frac{\kappa}{\sqrt q}.
\]
Similarly, the upper bound \(r<r_{\mathrm{gap}}^+=R_x^+=(a+b)\bigl(2m-a-b-\kappa/\sqrt q\bigr)\), divided by \(a+b>0\), gives
\[
\frac r{a+b}<2m-a-b-\frac{\kappa}{\sqrt q},
\]
hence, since \(c_z=2m\) and \(c_x/q=a+b+r/(a+b)\),
\[
c_z>\frac{c_x}{q}+\frac{\kappa}{\sqrt q}.
\]
By the spread-vertex certificate, \cref{obj:lem:spread-vertex-certificate}, the point
\[
\left(\frac{2m}{q},0,0\right)=\left(\frac{m}{m-1},0,0\right)
\]
belongs to \(T_m\) and is the strict unique minimizer of the reduced objective on \(T_m\).  % lean: spread_vertex_certificate

4.  The spread covariance is \(X(-1/(m-1),0)\), whose spectral coordinates are
\[
x=1-\left(-\frac1{m-1}\right)=\frac{m}{m-1}=\frac{2m}{q},\qquad
y=1-\frac{m-1}{m-1}=0,\qquad z_{\mathrm{sp}}=0.
\]
By \cref{obj:lem:block-spectral-coordinates}, the membership of this triple in \(T_m\) gives
\[
X_{\mathrm{spread}}\in\mathcal E_m^{\mathrm{blk}},
\]
and the value of \(F_{r,\kappa}\) at \(X_{\mathrm{spread}}\) is the reduced objective at
\((2m/q,0,0)\). If \(X\in\mathcal E_m^{\mathrm{blk}}\), write \(X=X(u,v)\). \Cref{obj:lem:block-spectral-coordinates} sends it to a point of \(T_m\) and identifies \(F_{r,\kappa}(X)\) with the reduced objective there. If this reduced point equals \((2m/q,0,0)\), then \(y=z_{\mathrm{sp}}=0\), so
\[
z_{\mathrm{sp}}-y=2mv=0
\quad\text{and}\quad
y+z_{\mathrm{sp}}=2+2(m-1)u=0,
\]
forcing \(v=0\) and \(u=-1/(m-1)\), hence \(X=X_{\mathrm{spread}}\). Therefore every \(X\ne X_{\mathrm{spread}}\) corresponds to a different point of \(T_m\), and the strict reduced optimality from Step 3 gives
\[
F_{r,\kappa}(X_{\mathrm{spread}})<F_{r,\kappa}(X).
\]  % lean: spread_certificate_relaxed_minimizer_and_gap

5.  Since \(m\) is odd, \cref{obj:lem:pm-reduced-slice-characterization} imposes the implementability condition
\[
y+z_{\mathrm{sp}}\ge d_m=\frac2m>0 .
\]
The spread point has \(y+z_{\mathrm{sp}}=0\), so
\[
X_{\mathrm{spread}}\notin\mathcal C_m^{\pm}.
\]
For the loss, pass from \((x,y,z_{\mathrm{sp}})\in T_m\) to the simplex variable
\[
t=(qx,\,y,\,z_{\mathrm{sp}}),
\qquad
t_x+t_y+t_z=qx+y+z_{\mathrm{sp}}=2m=M,
\]
a bijection between \(T_m\) and \(\Delta_M\) with inverse \((t_x,t_y,t_z)\mapsto(t_x/q,t_y,t_z)\).  With the coefficients and weights
\[
\alpha=\Bigl(\frac{c_x}{q},c_y,c_z\Bigr),
\qquad
\beta=\Bigl(\frac1q,1,1\Bigr),
\]
this bijection matches the two objectives exactly:
\[
\sum_i\alpha_it_i=\frac{c_x}{q}\,qx+c_yy+c_zz_{\mathrm{sp}}
=c_xx+c_yy+c_zz_{\mathrm{sp}},
\qquad
\sum_i\beta_it_i^2=\frac{(qx)^2}{q}+y^2+z_{\mathrm{sp}}^2=qx^2+y^2+z_{\mathrm{sp}}^2,
\]
so \(\Phi_{\alpha,\beta,\kappa}(t)=\phi_{r,\kappa}(x,y,z_{\mathrm{sp}})\), and the parity constraint reads \(t_y+t_z\ge d_m\).  Under this dictionary the relaxed minimizer of Step 3 is
\[
t^{\mathrm{rel}}=(2m,0,0),
\qquad
t^{\mathrm{rel}}_y+t^{\mathrm{rel}}_z=0<d_m ,
\]
so it is infeasible for the truncated simplex \(K_{d_m}\).  We check the data required by \cref{obj:lem:weighted-simplex-truncation}: the mass is \(M=2m>0\), the weights \(\beta=(1/q,1,1)\) are strictly positive with \(\beta_x=1/q=1/(2(m-1))\) and \(\beta_y=\beta_z=1\), the robustness weight satisfies \(\kappa\ge0\), and \(0<d_m=2/m\le1\le 2m=M\) since \(m\ge2\).  Its remaining input is the relaxed datum, namely the optimality of \(t^{\mathrm{rel}}\): transporting the strict uniqueness of Step 3 through the objective-matching bijection above, \(t^{\mathrm{rel}}\) is the strict unique minimizer of \(\Phi_{\alpha,\beta,\kappa}\) on \(\Delta_M\).  If \(\kappa>0\), \cref{obj:lem:weighted-simplex-active-set} supplies an admissible support--multiplier pair \((S,\lambda)\) whose active-set point minimizes \(\Phi_{\alpha,\beta,\kappa}\) over \(\Delta_M\), hence coincides with \(t^{\mathrm{rel}}\) by that uniqueness, so the truncation lemma may be applied to \((S,\lambda)\).  If \(\kappa=0\), the same lemma identifies the minimizers of the then-linear \(\Phi_{\alpha,\beta,\kappa}\) over \(\Delta_M\) with the exposed \(\alpha\)-minimizing face, so \(t^{\mathrm{rel}}\) is a point of that face; uniqueness makes the face the single point \(t^{\mathrm{rel}}\), which misses \(K_{d_m}\), so the tie convention attached to the \(\kappa=0\) branch is vacuous.  In either regime \(t^{\mathrm{rel}}_y+t^{\mathrm{rel}}_z=0<d_m\) puts us in the second alternative, which supplies a point \(t^\star\in K_{d_m}\) minimizing \(\Phi_{\alpha,\beta,\kappa}\) over \(K_{d_m}\); translating back, the corresponding \((x^\star,y^\star,z^\star_{\mathrm{sp}})\in T_m\) satisfies \(y^\star+z^\star_{\mathrm{sp}}\ge d_m\) and minimizes \(\phi_{r,\kappa}\) over the parity-truncated slice.  In particular both infima in \cref{obj:lem:rounding-gap-reduction} are attained:
\[
\inf_{(x,y,z_{\mathrm{sp}})\in T_m}\phi_{r,\kappa}
=\phi_{r,\kappa}\Bigl(\frac{2m}{q},0,0\Bigr),
\qquad
\inf_{\substack{(x,y,z_{\mathrm{sp}})\in T_m\\ y+z_{\mathrm{sp}}\ge d_m}}\phi_{r,\kappa}
=\phi_{r,\kappa}(x^\star,y^\star,z^\star_{\mathrm{sp}}).
\]
The truncated minimizer is feasible for the relaxed problem but differs from the relaxed minimizer, because \(y^\star+z^\star_{\mathrm{sp}}\ge d_m>0\) while the relaxed minimizer has \(y+z_{\mathrm{sp}}=0\); the strict uniqueness from Step 3 therefore gives
\[
\phi_{r,\kappa}\Bigl(\frac{2m}{q},0,0\Bigr)<\phi_{r,\kappa}(x^\star,y^\star,z^\star_{\mathrm{sp}}).
\]
Finally \cref{obj:lem:rounding-gap-reduction} identifies \(\Delta_m^{\pm}(r,\kappa)\) with the truncated reduced infimum minus the relaxed reduced infimum, so
\[
\Delta_m^{\pm}(r,\kappa)>0.
\]  % lean: spread_certificate_relaxed_minimizer_and_gap % lean: pos_gap_of_unique_min_outside_slice

6.  It remains to prove the even case. Assume \(m\) is even and let \(X\in\mathcal E_m^{\mathrm{blk}}\). Write \(X=X(u,v)\), and let \((x,y,z_{\mathrm{sp}})\) be its spectral coordinates. By \cref{obj:lem:block-spectral-coordinates}, this point lies in \(T_m\), so \(y,z_{\mathrm{sp}}\ge0\). For even \(m\), the parity threshold in \cref{obj:lem:pm-reduced-slice-characterization} is \(d_m=0\), and hence \(y+z_{\mathrm{sp}}\ge0\) automatically. The same lemma therefore gives \(X\in\mathcal C_m^{\pm}\). Thus
\[
\mathcal E_m^{\mathrm{blk}}\subseteq\mathcal C_m^{\pm}.
\]  % lean: blockElliptope_subset_implementable_of_even
\end{proof}
